% TOP_PROBLEM: {"id": 30007196, "title": "An infinite residually finite group with finite presentation and 2-power torsion", "category_id": 17, "category": {"id": 17, "name": "group_theory", "display_name": "Group Theory", "description": "Problems about groups, group actions, representations, and related algebraic structures.", "slug": "group-theory", "order_index": 17, "created_at": "2026-07-31T15:26:25.671Z"}, "status": "open"}
% FIRST_POSTED: 2026-10-06
% ATTEMPT_STATUS: solved
% OPENAI_CLAIM: {"schema_version":1,"source_repo":"https://github.com/openai/math","source_commit":"adc7f1241b42e322a6451854ab7e4b4c146bf78a","release_date":"2026-10-06","match":"full","resolution":"proved","independently_reviewed":false,"second_pass":"agreed","families":[{"family":"247","title":"An infinite finitely presented residually finite 2-group and a finitely presented nil algebra","subject":"Group theory","manuscripts":[{"title":"An infinite finitely presented residually finite 2-group","pdf_path":"preprints/An-infinite-finitely-presented-residually-finite-2-group-October-5-2026/residually-finite-torsion.pdf","readme_path":"preprints/An-infinite-finitely-presented-residually-finite-2-group-October-5-2026/README.md","date":"2026-10-05","theorem_ref":"preprints/An-infinite-finitely-presented-residually-finite-2-group-October-5-2026/build/sections/introduction.tex:12, result labelled thm:main (pinned adc7f1241b42e322a6451854ab7e4b4c146bf78a)"}],"lean":{"doc_path":"lean/docs/247.md","comparators":[],"declarations":[],"covers_main_theorem":false},"reasoning_trace":null}]}
\documentclass[11pt]{article}
\usepackage[margin=25mm]{geometry}
\usepackage{fontspec}
\setmainfont{Latin Modern Roman}
\usepackage{amsmath,amssymb,mathtools}
\usepackage{unicode-math}
\setmathfont{Latin Modern Math}
\usepackage{enumitem}
\usepackage{xurl,hyperref,bookmark}
\hypersetup{colorlinks=true,urlcolor=blue}
\setcounter{secnumdepth}{0}
\setlength{\parindent}{0pt}
\setlength{\parskip}{5pt}
\setlength{\emergencystretch}{3em}
\begin{document}
\section*{An infinite residually finite group with finite presentation and 2-power torsion}

\subsection{OpenAI's claim}
OpenAI presents a proof in the selected manuscript. The periodic group of the companion is claimed residually finite, with an infinite finite-index finitely presented subgroup whose element orders are finite powers of two and unbounded.

\subsection{Outline of the argument}
Phases of long graded words make finite-degree truncations detect group elements. Faithful Steinberg strips, cyclic-shift comparisons, and long-entry commutator calculations treat the central kernel as well, so residual finiteness survives passage to the periodic two-group subgroup.

\subsection{Scope relative to this problem}
The October main theorem asserts all four requested properties simultaneously: infinite, ordinarily finitely presented, residually finite and every element finite 2-power order. lean/docs/247.md covers the older generic periodic companion and explicitly does not document this stronger residual-finite 2-group main assertion. Within these assumptions, this is a claimed affirmative answer to the stated question.

\subsection{Formal verification}
For family 247, the selected entry gives an infinite ordinarily finitely presented periodic group. A uniform exponent, the nil-algebra conclusions, and the later residual-finiteness strengthening are not part of it. No Comparator is assigned to the selected manuscript here; the family page selects other papers or supporting statements. This site has not executed the Lean code. The scope descriptions do not certify the remaining manuscript arguments.

\subsection{Status caveat}
OpenAI attributes these manuscripts to work by an internal model. Its release cautions that arguments without formalization can contain mistakes. This catalogue has not independently checked the proofs; the labels report the mathematical claims attributed to OpenAI.

\subsection{Sources}
\begin{itemize}
\item OpenAI, \emph{An infinite finitely presented residually finite 2-group}, 2026-10-05; citation key \texttt{OAI:An-infinite-finitely-presented-residually-finite-2-group-October-5-2026}. \href{https://github.com/openai/math/blob/adc7f1241b42e322a6451854ab7e4b4c146bf78a/preprints/An-infinite-finitely-presented-residually-finite-2-group-October-5-2026/residually-finite-torsion.pdf}{PDF (GitHub)}; \href{https://github.com/openai/math/raw/adc7f1241b42e322a6451854ab7e4b4c146bf78a/preprints/An-infinite-finitely-presented-residually-finite-2-group-October-5-2026/residually-finite-torsion.pdf}{PDF (direct)}; \href{https://github.com/openai/math/blob/adc7f1241b42e322a6451854ab7e4b4c146bf78a/preprints/An-infinite-finitely-presented-residually-finite-2-group-October-5-2026/README.md}{Manuscript page and citation}; \href{https://github.com/openai/math/blob/main/preprints/An-infinite-finitely-presented-residually-finite-2-group-October-5-2026/residually-finite-torsion.pdf}{Latest version}.
\item Family 247: \href{https://github.com/openai/math/blob/adc7f1241b42e322a6451854ab7e4b4c146bf78a/lean/docs/247.md}{Lean scope}.
\item \href{https://github.com/openai/math/blob/adc7f1241b42e322a6451854ab7e4b4c146bf78a/CONTENTS.md}{OpenAI catalogue entry}.
\end{itemize}
\end{document}
